\documentclass[12pt]{article}
\usepackage[utf8]{inputenc}
\usepackage[T1]{fontenc}
\usepackage{hyperref}
\title{Patani Language Resources for K-12 Curriculum Development}
\author{Compiled by Claude Sonnet 5.5 (claude-sonnet-5-5) for Andr\'as Kornai}
\date{}
\begin{document}
\maketitle

\section{Introduction}
Patani (ISO 639-3: ptn; Glottocode pata1260) is an Austronesian language of the South Halmahera--West New Guinea subgroup, spoken in East Halmahera Regency, North Maluku, Indonesia. Published speaker counts vary (about 10,600 in Wikipedia, citing a 2000 figure; about 17,000 in Joshua Project, citing Ethnologue). Search found no school materials, primers or textbooks in Patani. What exists is mainly linguistic description. Indonesian is the language of trade and presumably of schooling.

\section{Primary Educational Resources}
No primers, textbooks, readers or school programmes in Patani were found in the searches carried out. Any K-12 curriculum would have to be built from scratch, drawing on the descriptive linguistic work listed below.

\section{Linguistic and Reference Materials}
\begin{itemize}
\item Wikipedia overview of the Patani language: \url{https://en.wikipedia.org/wiki/Patani_language}
\item Linn Iren Sj{\aa}nes R{\o}dvand (2024), a grammatical description of Patani (University of Oslo), as listed in Glottolog: \url{https://glottolog.org/resource/reference/id/704611}
\item Glottolog language page (Glottocode pata1260): \url{https://glottolog.org/resource/languoid/id/pata1260.bigmap.html}
\item Study of the possessive system of Patani (De Gruyter, 2023): \url{https://www.degruyterbrill.com/document/doi/10.1515/stuf-2023-2016/html?lang=de}
\item ASJP word list for Patani: \url{https://asjp.clld.org/languages/PATANI}
\item Comparative study of Patani and East Makeang (IJTK, IAIN Kendari): \url{https://ijtk.iainkendari.ac.id/index.php/IJTK/article/download/33/36/92}
\end{itemize}

\section{Religious and Community Resources}
\begin{itemize}
\item Joshua Project profile of the Patani people in Indonesia (demographic and scripture-status information): \url{https://joshuaproject.net/people_groups/14335}
\item Asia Harvest people-group profile of the Patani: \url{https://people-groups.asiaharvest.org/Indonesia/Patani.pdf}
\end{itemize}
No Patani scripture translation or community organisation was identified in the search results.

\section{Digital and Audio Resources}
No audio, video, corpus or software resources specific to Patani were found in the searches. The ASJP word list above is the only lexical data located online.

\section{Curriculum Development Context}
The language is spoken by all adults but not all young people, and it is not known to be taught in schools. The evidence found supports a documentation-first approach: the Rødvand grammar and the possession study could inform an orthography and teaching grammar, while primers, vocabulary lists and recordings would need to be produced with the community. Wikipedia and Joshua Project figures differ, so the speaker population should be verified locally.

\end{document}